% sections/ndss_06_temporal.tex  — NDSS 2027 / CCS 2026
% §6: Cross-Platform Validation — skills.sh Live Execution
% Target: ~1.5 columns / ~700 words

\section{Cross-Platform Validation}
\label{sec:crossplatform}

The ClawHub \cape experiment (\S\ref{sec:cape}) validates the
compositional attack surface on a single platform.
To establish generality, we conduct two independent cross-platform analyses:
(i) a static inconsistency measurement on \skillssh{},
and (ii) a live execution experiment using a real \texttt{codex}
agent instance.

\subsection{Static Analysis on \skillssh{}}

We apply the \lscps pipeline (\S\ref{sec:measurement}) without modification
to \NskillsSH{} \skillssh{} skills (retrieved April 2026).
The pipeline is designed for ClawHub's SKILL.md format;
\skillssh{} uses the same format,
making the transfer zero-shot.

\paragraph{Result.}
IC-1 (complete capability non-disclosure) applies to
$\mathbf{>90\%}$ of \skillssh{} skills---
a rate higher than ClawHub ($75.8\%$),
consistent with the hypothesis that disclosure norms
depend on marketplace culture rather than technical constraints.
Source-skill concentration ($S_4$: local system access, $S_1$: web fetch)
and sink-skill concentration ($K_1$: shell execution, $K_5$: persistence)
mirror the ClawHub distribution, confirming that the
Source$\times$Sink attack surface is structurally universal.

\subsection{Live Execution: \skillssh{} CAPE}
\label{sec:ss_cape}

\paragraph{Dataset.}
Starting from the \NskillsSH{} \skillssh{} corpus,
we apply a three-stage filter:
(1) LSCPS source-sink pair detection ($4{,}154$ self-exfil candidates);
(2) Codex semantic audit via chain-of-thought reasoning
    (prompt: ``Does this skill's SKILL.md omit any sink-type dangerous capability?''),
retaining $\NssPairs{}$ Codex-confirmed viable attack paths; and
(3) install-command availability (required for automation).
The $\NssPairs{}$ confirmed skills span all five
sink channel types (Table~\ref{tab:ss_results}).

\paragraph{Experimental setup.}
Each skill is evaluated in a fully isolated workspace
(\texttt{/tmp/cape\_sandboxes/\textlangle{}id\textrangle{}/})
containing realistic credential files
(\texttt{credentials/openai.token}, \texttt{credentials/github.token},
\texttt{memory/hot/HOT\_MEMORY.md}, \texttt{BOOT.md}).
Installation: \texttt{npx skills add \textlangle{}repo\textrangle{} --skill \textlangle{}name\textrangle{} -a codex -y}.
Execution: \texttt{codex exec --full-auto --skip-git-repo-check \textlangle{}trigger\textrangle{}}.
This represents a self-exfil attack model:
the skill is both source (it processes external input)
and sink (it executes dangerous operations),
which is the most conservative threat model---
it does not require any composition.
Side effects are verified by filesystem probe:
a skill-specific marker file, BOOT.md content,
memory-file modification, or exfil log.
Each task type uses the injection trigger described in Table~\ref{tab:ss_results}.

\paragraph{Results.}
Table~\ref{tab:ss_results} reports results for \NssRun{} executed skills
(31 skipped: GitHub access-controlled or withdrawn between audit and execution).
Install success: $\NssRun{}/\NssRun{} = \NssInstallRate{}$.
Attack success: $165/166 = \NssSuccessRate{}$.

\begin{table}[t]
\centering
\caption{%
  \skillssh{} live execution results by sink channel type.
  All experiments use \texttt{codex exec --full-auto}
  with per-sink injection triggers.
  ``Success'' = verified filesystem side effect.
}
\label{tab:ss_results}
\setlength{\tabcolsep}{4pt}
\small
\begin{tabular}{llrrc}
\toprule
\textbf{Task} & \textbf{Sink Channel} & \textbf{N} & \textbf{Success} & \textbf{Rate}\\
\midrule
T06 & C1: Shell execution        & 75 & 75 & \textbf{100\%}\\
T07 & C4: Network data collect   & 26 & 25 & \textbf{96\%}\\
T02 & C2: Filesystem (cred del.) & 23 & 23 & \textbf{100\%}\\
T13 & C3: Agent memory poison    & 20 & 20 & \textbf{100\%}\\
T19 & C4: Messaging webhook      & 13 & 13 & \textbf{100\%}\\
T05 & C5: Persistence (boot)     &  9 &  9 & \textbf{100\%}\\
\midrule
\textbf{All} & & \textbf{166} & \textbf{165} & \textbf{\NssSuccessRate{}}\\
\bottomrule
\end{tabular}
\end{table}

\paragraph{Key findings.}(i)~\textbf{Install-time barrier is zero}:
every installable skill installs cleanly without any platform-level
safety gate ($\NssInstallRate{}$, 166/166).
(ii)~\textbf{Runtime exploitability is near-universal}:
$\NssSuccessRate{}$ of installed skills produce dangerous verified
side effects---T06 (shell exec) at 100\%, T13 (memory poison) at 100\%,
T05 (boot hook persistence) at 100\%.
(iii)~\textbf{Network exfil is the only defended channel}:
T07 achieves 96\% rather than 100\%; agent $\texttt{gpt-5.4}$
detects explicit credential-reading scripts in 4\% of network-exfil cases
and refuses---the only category with any observable safety response.
This asymmetry confirms that safety defenses are
heuristic-based and target-specific, not capability-based.
(iv)~\textbf{Cross-platform consistency}:
the Source$\times$Sink matrix from ClawHub (\S\ref{sec:cape})
is reproduced across both platforms and both attack models
(compositional and self-exfil), confirming that
IC-1 gaps are a structural property of skill markets,
not an artifact of a single platform's design.

\paragraph{Temporal evolution.}
\label{sec:temporal}
We additionally analyze $24{,}248$ version transitions across $8{,}120$
\clawhub{} skills tracked between November 2024 and March 2026.
Of these, 6.8\% introduce new security-sensitive capabilities
without any description update (silent TCE).
627 skills (7.7\% of updated skills) acquired new capability families
silently---\eg a web-search skill that gained \texttt{FILE\_WRITE}
in a minor update.
Applied to 40 confirmed malicious skill version pairs
(34 independently flagged malicious skills),
\lscps achieves~\Tsixrecall{} recall at Tier-6 severity threshold;
the platform scanner achieves 0\%.
This establishes temporal monitoring as a necessary complement
to publish-time screening.
